% TOP_PROBLEM: {"id": 3400, "title": "Linear-size circuits for stable $0,1 < 2$ sorting?", "category_id": 15, "category": {"id": 15, "name": "computer_science", "display_name": "Computer Science", "description": "Computational complexity, algorithms, and theoretical CS.", "slug": "computer-science", "order_index": 15, "created_at": "2026-07-31T15:26:25.671Z"}, "status": "open", "problem_number": "OPG-474", "set_id": 12}
% FIRST_POSTED: 2026-10-02
% ATTEMPT_STATUS: unresolved
% UnsolvedMath ID 3400; source catalogue code OPG-474
% !TeX program = lualatex
% Research attempt by GPT-6 Astra Ultra; no claim of a new complete solution.
\documentclass[11pt,a4paper]{article}
\usepackage[margin=25mm]{geometry}
\usepackage{fontspec}
\setmainfont{lmroman10-regular.otf}[BoldFont=lmroman10-bold.otf,
 ItalicFont=lmroman10-italic.otf,BoldItalicFont=lmroman10-bolditalic.otf]
\usepackage{amsmath,amssymb,mathtools}
\usepackage{unicode-math}
\setmathfont{latinmodern-math.otf}
\AtBeginDocument{\let\setminus\smallsetminus}
\usepackage{xurl,hyperref}
\hypersetup{colorlinks=true,urlcolor=blue}
\setcounter{secnumdepth}{0}
\setlength{\parindent}{0pt}
\setlength{\parskip}{5pt}
\setlength{\emergencystretch}{3em}
\begin{document}
\section*{Linear-size circuits for stable $0,1 < 2$ sorting?}\label{3400}
{\small UnsolvedMath ID 3400 · OPG-474 · Computer Science · OpenGarden}
\subsection{Definitions and mathematical statement}
For a word $x=x_1\cdots x_n\in\{0,1,2\}^n$, let $u(x)$ be the
subword obtained by deleting every occurrence of $2$, preserving
the order of all remaining symbols, and let $t(x)$ count the deleted
symbols. Define $f_n(x)=u(x)2^{t(x)}$, where the superscript here
denotes repetition of a symbol. Thus $f_n$ moves all twos to the
right without changing the relative order of the zeros and ones.

Encode the three symbols by $00,01,10$, respectively. The question
is whether there is a constant $C$ and, for every $n\ge1$, a Boolean
circuit on $2n$ input and $2n$ output bits computing $f_n$ on valid
encodings with at most $Cn$ gates. Use fan-in-two AND and OR gates
and fan-in-one NOT gates, with unrestricted reuse of outputs.
The circuit's behavior on invalid symbol encodings is immaterial.

This is a nonuniform circuit-size question; no bound on depth or
procedure for constructing the circuits is required. The choice of
a fixed finite complete bounded-fan-in basis affects size by only
a constant factor. Stability is essential: sorting all zeros before
all ones would discard the original relative order and compute a
different function. Determine whether these stable-compaction maps
have linear circuit size.
\subsection{Short English statement}
Can bounded-fan-in Boolean circuits of linear size move all twos to the end of a ternary word while preserving the order of its zeros and ones?
\subsection{Research attempt}
\paragraph{Scope and process.}
This is a GPT-6 Astra Ultra research attempt dated 2 October 2026, using
primary-source browsing and elementary mathematical checks. The canonical
definition above is retained. Results below are restricted results or
reductions, not a claim of a new solution to the entire catalogue question.
No formal proof assistant or external mathematical referee checked this text.


\paragraph{Literature and exact task.}
The source poses a Boolean-circuit question about retaining the order
of zeroes and ones while moving twos to the end. Ordinary sorting
by the three numerical values changes that order and computes a
different function. Asharov, Lin and Shi [R1] obtain strong circuit
bounds for short-key sorting, but explicitly describe their compaction
construction as non-stable. We do not transfer that construction to
this target. Here we give an explicit quadratic-size stable circuit
and an elementary linear lower bound in the required gate model.

\paragraph{One adjacent operation with a constant gate cost.}
For adjacent symbols $a,b$, exchange them exactly when $a=2$ and
$b\ne2$. Otherwise retain them. This operation never exchanges a zero
with a one. Using the prescribed encodings, let the two input pairs be
$(a_1,a_0)$ and $(b_1,b_0)$. On valid encodings, the exchange flag is
\[
                         s=a_1\wedge\neg b_1.
\]
For each of the four output bits, let $u$ be its bit if no exchange
occurs and $v$ its bit if exchange occurs. The output is the multiplexer
\[
                      (\neg s\wedge u)\vee(s\wedge v).
\]
Computing $\neg b_1$, $s$, and $\neg s$ uses three gates, and the four
multiplexers use three gates each. Thus fifteen gates suffice per
adjacent operation. Fan-out shares $s$ and $\neg s$ freely. Behaviour
on the unused encoding $11$ is irrelevant, as permitted by the
statement, while valid input pairs always produce valid outputs.

\paragraph{A fixed network computing the required function.}
First apply the adjacent operation successively at positions
$(1,2),(2,3),\ldots,(n-1,n)$. Then repeat on the prefix of length
$n-1$, and continue down to a prefix of length two. The locations are
fixed before seeing the input, so this is a circuit, not an algorithm
with input-dependent wiring. There are
\[
                       \sum_{j=1}^{n-1}j=\frac{n(n-1)}2
\]
operations. During each left-to-right pass, a two moves to the last
position of the prefix if that prefix contains any two. If it contains
none, it already has the desired form and the pass makes no change.
Induction on prefix length consequently puts every two after every
non-two. At every stage the multiset of symbols is preserved, and no
operation exchanges two non-twos. Their relative order is therefore
unchanged. The final word is exactly $u(x)2^{n-|u(x)|}$.
The resulting circuit has size at most $15n(n-1)/2$ for $n\ge2$;
for $n=1$ direct wires suffice. This proves an explicit $O(n^2)$ upper
bound including stability, with no arithmetic or addressing primitive
hidden in the gate count.

\paragraph{An elementary lower bound.}
Restrict inputs to symbols zero and two. Their high encoding bits
$z_1,\ldots,z_n$ are independent Boolean variables and their low bits
are zero. The first output symbol is two precisely when every input
symbol is two. Hence its high bit computes
\[
                           z_1\wedge\cdots\wedge z_n.
\]
This function depends essentially on all $n$ variables. Consider the
ancestor subgraph of that output in any circuit after the restriction.
Let $b$ and $u$ count its binary and unary gates. It contains at least
$n$ variable-source vertices and is connected when directions are
forgotten. If it has additional constant sources, they only increase
the required edge count. Thus it has at least $n+b+u-1$ edges, while
its gate fan-in supplies at most $2b+u$ edges. It follows that
$b\ge n-1$. Therefore any circuit for the full function needs at
least $n-1$ gates. This argument allows arbitrary fan-out and makes
no comparison-only assumption.

\paragraph{Verification and the gap in a linear scan.}
The standard-library certificate [R2] checks the fifteen-gate operation
on all nine valid symbol pairs. It then exhausts all ternary words
through length eight and compares the fixed network with directly
deleting twos and appending them. These are functional checks, not a
proof of an asymptotic lower bound; the induction and dependency count
provide the general results.

A sequential linear-time scan can append each non-two to a growing
output list, but implementing that description requires a
variable-address write or equivalent routing. Such a primitive is
not one AND, OR or NOT gate. Counting each append as one operation
therefore does not produce an $O(n)$ Boolean circuit. Conversely,
lower bounds for comparison networks alone would not rule out the
arbitrary Boolean circuits allowed here.

\paragraph{Final status and first remaining gap.}
\textbf{UNRESOLVED.} An explicit stable quadratic circuit and a linear
lower bound are verified. The remaining question is whether the stable
routing can be accomplished with only a constant number of Boolean
gates per input symbol. Neither the adjacent network nor the
non-stable short-key sorting result establishes that bound.

\subsection{Sources}
\begin{itemize}
\item[S1] ULAM.AI and the UnsolvedMath Contributors, supplied problems.json, record 3400 (OPG-474), OpenGarden collection. Catalogue identity and source transcription; CC BY 4.0.
\url{https://huggingface.co/datasets/ulamai/UnsolvedMath}.
\item[S2] Open Problem Garden, Linear-size circuits for stable $0,1 < 2$ sorting?. Original problem and discussion; the mathematical formulation below is expanded from the supplied source record.
\url{http://www.openproblemgarden.org/op/linear_size_circuits_for_stable_0_1_2_sorting}.

\item[R1] Gilad Asharov, Wei-Kai Lin and Elaine Shi, \emph{Sorting
Short Keys in Circuits of Size $o(n\log n)$}, 2020 preprint; SIAM
Journal on Computing 51 (2022), 424--466. Primary PDF, including the
stability discussion in its introduction, checked 2 October 2026.
\url{https://arxiv.org/abs/2010.09884}.
\item[R2] Reproducible exhaustive checks:
\path{scripts/check_attempt_batch03_colouring_2026_10_02.py}, section 3400.

\end{itemize}
\end{document}
